
\medskip
\noindent \SGARefTarget{sga7.sga7i.eviii.n4-12}{4.12}.  Let $P$ be a smooth connected group scheme, let
$\mathbb L$ be a line bundle on $P$, and put
\begin{equation}\tag{\SGARefTarget{sga7.sga7i.eviii.n4-12-1}{4.12.1}}
 \delta(\mathbb L)=
 \pi^*\mathbb L\otimes
 \pr_1^*(\mathbb L)^{-1}\otimes
 \pr_2^*(\mathbb L)^{-1},
\end{equation}
with the notation of \SGARefLink{sga7.sga7i.eviii.n4-2}{4.2}.  As already observed in VII \SGARefLink{sga7.sga7i.evii.n1-3-4}{1.3.4}, this is
a birigidified line bundle on $P\times P$.  By \SGARefLink{sga7.sga7i.eviii.n4-3}{4.3}, if
$\delta(\mathbb L)$ comes from a biextension, that biextension is
determined up to unique isomorphism.  In particular, by \SGARefLink{sga7.sga7i.eviii.n4-7}{4.7}, it is
defined when $P$ is an extension of an abelian scheme $A$ by an
iterated extension of tori and copies of $\underline G_a$.  Whenever
it is defined, this biextension is plainly symmetric.

In the preceding favourable case, and assuming that the maximal
torus of $P$ is split, the construction of $\delta(\mathbb L)$ can
be reduced to the case where $P$ is an abelian scheme.  Indeed, the
construction is functorial in $P$, while the homomorphism
\begin{equation}\tag{\SGARefTarget{sga7.sga7i.eviii.n4-12-2}{4.12.2}}
 \Pic(A)\longrightarrow\Pic(P)
\end{equation}
induced by the canonical quotient homomorphism
$f:P\longrightarrow A$ is surjective by \SGARefLink{sga7.sga7i.eviii.n4-5}{4.5} (cf. the beginning of
the proof of \SGARefLink{sga7.sga7i.eviii.n4-6}{4.6}).  We can therefore find an isomorphism
\[
 \mathbb L\simeq f^*\mathbb M,
\]
where $\mathbb M$ is a line bundle on $A$.  Then
$\delta(\mathbb L)$ is the pullback of the biextension
$\delta(\mathbb M)$ of $(A,A)$ by $\underline G_{m,k}$.

\medskip
\noindent\underline{Proposition \SGARefTarget{sga7.sga7i.eviii.n4-13}{4.13}.}  \underline{Under the
hypotheses of \SGARefLink{sga7.sga7i.eviii.n4-12}{4.12}, suppose that $\mathbb L$ is an ample line bundle
on $P$.  If the biextension $\delta(\mathbb L)$ of $(P,P)$ by
$\underline G_{m,k}$ is defined, then it is nondegenerate in the
sense of \SGARefLink{sga7.sga7i.eviii.n4-11}{4.11}.  If $P$ is an extension of an abelian scheme $A$ by a
smooth connected affine group $L$ satisfying the condition of \SGARefLink{sga7.sga7i.eviii.n4-7}{4.7},
then the biextension of $(A,A)$ from which $\delta(\mathbb L)$ comes
by \SGARefLink{sga7.sga7i.eviii.n4-8}{4.8} is of the form $\delta(\mathbb M)$, where $\mathbb M$ is an
ample line bundle on $A$.}

We reduce to the case where $k$ is algebraically closed.  By the
definition, it is then enough, with the notation of \SGARefLink{sga7.sga7i.eviii.n4-12}{4.12}, to prove
that the biextension $\delta(\mathbb M)$ of $(A,A)$ is
nondegenerate.  M. Raynaud proves in his thesis [2, XI 1.11] that if
$\mathbb L=f^*\mathbb M$ is ample, then $\mathbb M$ is ample as
well.  The theory of A. Weil then gives the familiar conclusion that
$\delta(\mathbb M)$ is nondegenerate [1, 2], that is, that the
homomorphism from $A$ to its dual abelian scheme is an isogeny.

\medskip
\noindent\underline{Remark \SGARefTarget{sga7.sga7i.eviii.n4-14}{4.14}.}  Suppose we are in the favourable
case where $P$ is an extension of an abelian scheme $A$ by a group
$L$ which is an iterated extension of tori and copies of
$\underline G_a$, and consider the canonical composite
\begin{equation}\tag{\SGARefTarget{sga7.sga7i.eviii.n4-14-1}{4.14.1}}
 \operatorname{NS}(A)
 \xrightarrow{\ \delta\ }
 \Biext^1(A,A;\underline G_m)
 \xrightarrow{\ \sim\ }
 \Biext^1(P,P;\underline G_m).
\end{equation}
By \SGARefLink{sga7.sga7i.eviii.n4-8}{4.8} the second arrow is bijective.  On the other hand, the theory
of abelian varieties shows that the first arrow identifies the
Neron--Severi group $\operatorname{NS}(A)$ of $A$ with the subgroup
of
\[
 \Biext^1(A,A;\underline G_m)=\operatorname{Corr}(A,A)
\]
consisting of the classes of symmetric biextensions, or equivalently
symmetric divisorial correspondences.  Thus (\SGARefLink{sga7.sga7i.eviii.n4-14-1}{4.14.1}) identifies
$\operatorname{NS}(A)$ with the group of symmetric biextensions of
$(P,P)$ by $\underline G_m$.  By \SGARefLink{sga7.sga7i.eviii.n4-7}{4.7} this group can equally be
described as the group of classes of symmetric divisorial
correspondences on $(P,P)$.

It is natural to call a divisorial correspondence on $(P,P)$
\underline{ample}, and likewise to call the corresponding
biextension of $(P,P)$ by $\underline G_m$ \underline{ample}, when
it comes from an ample element of $\operatorname{NS}(A)$.  This
condition is plainly invariant under extension of the base field;
the notion therefore still makes sense without the preceding
restrictive hypothesis on $P$, assuming only that $P$ is smooth and
connected.

Suppose further that $L$ is ``split solvable over $k$'', in the
sense that it admits a devissage by copies of $\underline G_a$ and
$\underline G_m$.  The preceding observation and the surjectivity of
(\SGARefLink{sga7.sga7i.eviii.n4-12-2}{4.12.2}) show that the image of
\begin{equation}\tag{\SGARefTarget{sga7.sga7i.eviii.n4-14-2}{4.14.2}}
 \Pic(P)\xrightarrow{\ \delta\ }
 \Biext^1(P,P;\underline G_m)
\end{equation}
consists of the symmetric biextensions, and that its kernel is the
subgroup $\Pic^0(P)$ of $\Pic(P)$ which is the image of $\Pic^0(A)$.
Consequently, if we put
\[
 \operatorname{NS}(P)=\Pic(P)/\Pic^0(P)
 \simeq\operatorname{NS}(A),
\]
we may also interpret (\SGARefLink{sga7.sga7i.eviii.n4-14-1}{4.14.1}) as a canonical isomorphism
\begin{equation}\tag{\SGARefTarget{sga7.sga7i.eviii.n4-14-3}{4.14.3}}
 \operatorname{NS}(P)
 \xrightarrow{\ \sim\ }
 \Biext^1(P,P;\underline G_m)^{\mathrm{sym}}.
\end{equation}

Finally, the theorem of Raynaud cited in \SGARefLink{sga7.sga7i.eviii.n4-13}{4.13} says that an element
$\xi$ of $\Pic(P)$ is ample if and only if its image
$\delta(\xi)$ under (\SGARefLink{sga7.sga7i.eviii.n4-14-2}{4.14.2}) is the class of an ample biextension.
